\section{Introduction}
\label{sec:intro}

In modern machine learning pipelines, the rapid growth of specialized, task-specific deep neural networks has created a severe deployment bottleneck. For practitioners operating under real-world constraints—such as limited hardware memory, strict latency budgets, and reduced edge storage—serving multiple specialized expert models is commercially and technically unfeasible~\cite{gholami2021survey}. Each additional downstream task requires storing and serving an independent, full-sized checkpoint, causing aggregate system storage and inference-time memory to scale linearly with the number of tasks. 

To mitigate this burden, weight-space \emph{model merging} has emerged as an attractive, zero-shot paradigm for multi-task fusion~\cite{ilharco2022editing, wortsman2022model}. By exploiting the linear mode connectivity of deep networks fine-tuned from a shared pre-trained checkpoint~\cite{frankle2018lottery, neyshabur2020shared}, techniques such as Task Arithmetic (TA)~\cite{ilharco2022editing} extract task-specific parameter updates (task vectors) and average them into a single joint checkpoint. This process enables a single model to concurrently serve multiple downstream tasks with zero retraining cost and zero increase in inference-time parameters.

However, in practical deployment scenarios, weight fusion is only half the battle. To run on commodity edge hardware, specialized embedded processors, or low-latency cloud servers, models must undergo \emph{Post-Training Quantization} (PTQ) to low-bit formats, such as 8-bit (INT8) or 4-bit (INT4) representations~\cite{jacob2018quantization, gholami2021survey}. PTQ is the standard industry practice because it reduces the memory footprint by up to $4\times$, slashes memory bandwidth bottlenecks, and activates high-throughput integer tensor cores. 

\begin{figure}[t]
  \centering
  \includegraphics[width=\columnwidth]{qmerge_vs_baselines.png}
  \caption{Comparison of average multi-task accuracy (\%) across three independent seeds for 8-bit and 4-bit quantization configurations. Under 8-bit PTQ, our proposed Q-Merge (Adam GD with STE) not only recovers the quantization loss but actively exceeds both the unquantized FP16 baseline and the true unquantized optimized AdaMerging ceiling. Under 4-bit PTQ (using standard per-channel quantization), models no longer collapse; our Q-Merge (Adam GD) achieves a substantial \textbf{63.36\%} average accuracy, outperforming the naive baseline by \textbf{6.70\%} absolute.}
  \label{fig:qmerge_vs_baselines}
\end{figure}

Crucially, standard model merging research has largely ignored these severe deployment compression constraints. When weight merging and post-training quantization collide, standard deployment pipelines suffer from a fundamental mismatch:
\begin{enumerate}
    \item \textbf{Merge-then-Quantize (M-then-Q):} Standard Task Arithmetic merges full-precision models using static, uniform coefficients and then quantizes the joint checkpoint. However, the uniform fusion of weight coordinates creates delicate, multi-task decision boundaries that are highly sensitive to coordinate-wise noise. Under PTQ rounding, quantization noise scrambles these fine-tuned boundaries, resulting in significant multi-task accuracy degradation.
    \item \textbf{Quantize-then-Merge (Q-then-M):} Quantizing each task-specific expert to low-bit formats first and then merging their discrete integer weights completely breaks linear mode connectivity. The discrete, non-overlapping quantization grids of independent experts cannot be algebraically aligned, yielding a non-functional joint model.
\end{enumerate}

To bridge this gap and provide an efficient, robust solution for edge deployment, we propose \textbf{Quantization-Aware Model Merging (Q-Merge)}. Guided by a pragmatic research philosophy that prioritizes real-world utility, simplicity, and low computational overhead, Q-Merge optimizes layer-wise merging coefficients \emph{directly under the quantization operator}. Q-Merge operates at test-time using a tiny, unlabeled calibration stream (just 64 images in total) to minimize the joint classification entropy. This makes the method data-private, calibration-free, and extremely fast, requiring no expensive training loops.

In this work, we conduct a rigorous empirical study of Q-Merge across two optimization paradigms:
\begin{itemize}
    \item \textbf{Zero-Order Optimization (1+1 ES):} A derivative-free black-box mutation optimizer that stochastically updates merging coefficients, treating the quantized model as a non-differentiable oracle~\cite{yang2024adamerging}.
    \item \textbf{First-Order Optimization (Adam GD with STE):} Gradient-based optimization utilizing the Straight-Through Estimator (STE)~\cite{bengio2013estimating} to propagate gradients from the entropy loss back through the non-differentiable rounding operator to the continuous merging coefficients.
\end{itemize}

Through extensive evaluation across three independent random seeds on a multi-task vision benchmark (MNIST, FashionMNIST, CIFAR-10, SVHN) using a pre-trained ViT-Tiny backbone (as illustrated in Figure~\ref{fig:qmerge_vs_baselines}), we report the following key contributions:
\begin{itemize}
    \item \textbf{Overcoming the 8-Bit Quantization Gap:} We demonstrate that naive post-merge quantization (M-then-Q) degrades multi-task accuracy. Q-Merge using Adam GD with STE actively compensates for this noise, achieving an average multi-task accuracy of \textbf{74.30\%}, which remarkably surpasses the unquantized FP16 baseline (\textbf{71.88\%}) and standard AdaMerging (\textbf{73.21\%}), while recovering $99.9\%$ of the true unquantized Adam-optimized ceiling (\textbf{74.38\%}). We thoroughly deconstruct the optimizer confounding factor to show that this boost is enabled by unlocking highly efficient first-order gradient flow through the non-differentiable quantization operator using the STE.
    \item \textbf{First-Order vs. Zero-Order Stability:} We show that direct gradient feedback enabled by STE is highly superior to zero-order random-walk mutations. Q-Merge (Adam GD) converges faster, achieves higher final performance, and exhibits over $2.7\times$ lower seed-to-seed variance, confirming that first-order optimization is highly reliable even in rugged, quantized coordinate spaces.
    \item \textbf{Unlocking 4-Bit Model Merging:} We demonstrate that under aggressive 4-bit quantization, naive post-training rounding on a per-tensor level collapses the model. However, by employing standard per-channel (channel-wise) weight quantization, linear connectivity is preserved. Our proposed Q-Merge (Adam GD with STE) further optimizes this alignment to achieve an outstanding average accuracy of \textbf{63.36\%}, outperforming the naive post-merge quantization baseline (\textbf{56.66\%}) by \textbf{6.70\%} absolute and post-hoc quantized AdaMerging (Adam GD) (\textbf{62.01\%}) by \textbf{1.35\%} absolute, proving that extreme low-bit model merging is highly viable and practical when combined with quantization-aware optimization.
\end{itemize}